\documentclass[12pt,a4paper]{report}
\usepackage[utf8]{inputenc}
\usepackage[T1]{fontenc}
\usepackage[french]{babel}
\usepackage{amsmath}
\usepackage{amsfonts}
\usepackage{amssymb}
\usepackage{graphicx}
\usepackage[left=2cm,right=2cm,top=2.5cm,bottom=2.5cm]{geometry}
\usepackage{fancyhdr}
\usepackage{hyperref}
\usepackage{csquotes}
\usepackage{enumitem}
\usepackage{float}
\usepackage{hyperref}
\usepackage{booktabs}
\usepackage{array}
\usepackage{bookmark}
\usepackage{hyperref}
\usepackage{lipsum}

\hypersetup{
    colorlinks=true,
    linkcolor=black,
    filecolor=magenta,
    urlcolor=cyan,
    pdftitle={Rapport de Stage - AI Network Copilot},
    pdfpagemode=FullScreen,
    }

\urlstyle{same}

\pagestyle{fancy}
\fancyhf{}
\rhead{\small\emph{AI Network Copilot}}
\lhead{\small\emph{Rapport de Stage}}
\rfoot{\thepage}

\begin{document}

\begin{titlepage}
    \centering
    \vspace*{2cm}

    {\LARGE \textbf{AI Network Copilot}}\\[1.5cm]

    {\Large Rapport de Stage}\\[0.5cm]

    \vspace{1cm}

    {\large Réalisé par :}\\
    {\large [Votre Nom]}\\[0.5cm]

    {\large Encadrant :}\\
    {\large [Nom de l'encadrant]}\\[0.5cm]

    {\large Établissement :}\\
    {\large [Nom de l'établissement]}\\[0.5cm]

    {\large Département :}\\
    {\large [Nom du département]}\\[2cm]

    {\large Période : Juin - Juillet 2026}\\[1cm]

    {\large Version : 1.0}\\[0.5cm]

    {\large Date : \today}\\[2cm]

    \vfill

    \vfill
\end{titlepage}

\chapter*{Résumé}
\addcontentsline{toc}{chapter}{Résumé}
\thispagestyle{empty}

Ce rapport présente le projet \textit{AI Network Copilot}, réalisé dans le cadre d'un stage de quatre semaines. L'objectif du projet était de développer un outil de surveillance réseau intelligent capable de détecter des anomalies en temps réel grâce à l'apprentissage automatique.

L'architecture proposée se compose d'un backend développé avec FastAPI, responsable de la collecte des métriques réseau, de l'analyse des données et de la fourniture d'une API RESTful. Le frontend, développé avec Streamlit, offre un tableau de bord interactif permettant de visualiser les métriques, les alertes et la topologie du réseau en temps réel.

Le projet intègre un moteur de détection d'anomalies basé sur l'apprentissage automatique, qui analyse en continu les métriques telles que la latence, le taux d'erreur et le trafic pour identifier les comportements anormaux. Les résultats sont présentés sous forme de visualisations interactives, incluant des graphiques temporels, des cartes de répartition des alertes et une topologie réseau dynamique.

Les tests effectués ont démontré la capacité du système à détecter diverses anomalies réseau, telles que des latences élevées, des taux d'erreur anormaux et des pics de trafic. L'interface utilisateur permet une surveillance proactive et une réaction rapide aux incidents détectés.

En conclusion, le projet \textit{AI Network Copilot} constitue une solution prometteuse pour la surveillance réseau proactive, combinant des technologies modernes de développement web et d'apprentissage automatique pour offrir une visibilité en temps réel sur la santé du réseau.

\textbf{Mots-clés :} Surveillance réseau, Détection d'anomalies, Apprentissage automatique, FastAPI, Streamlit, WebSocket.

\chapter*{Abstract}
\addcontentsline{toc}{chapter}{Abstract}
\thispagestyle{empty}

This report presents the \textit{AI Network Copilot} project, carried out as part of a four-week internship. The project aimed to develop an intelligent network monitoring tool capable of detecting anomalies in real-time using machine learning.

The proposed architecture consists of a backend developed with FastAPI, responsible for collecting network metrics, analyzing data, and providing a RESTful API. The frontend, developed with Streamlit, offers an interactive dashboard for visualizing metrics, alerts, and network topology in real-time.

The project integrates a machine learning-based anomaly detection engine that continuously analyzes metrics such as latency, error rate, and traffic to identify abnormal behavior. Results are presented as interactive visualizations, including time-series charts, alert distribution maps, and dynamic network topology.

Tests conducted have demonstrated the system's ability to detect various network anomalies, such as high latency, abnormal error rates, and traffic spikes. The user interface enables proactive monitoring and rapid response to detected incidents.

In conclusion, the \textit{AI Network Copilot} project constitutes a promising solution for proactive network monitoring, combining modern web development technologies and machine learning to provide real-time visibility into network health.

\textbf{Keywords:} Network monitoring, Anomaly detection, Machine learning, FastAPI, Streamlit, WebSocket.

\tableofcontents
\thispagestyle{empty}

\chapter{Remerciements}
\addcontentsline{toc}{chapter}{Remerciements}
\thispagestyle{empty}

Je souhaite exprimer ma gratitude envers toutes les personnes qui ont contribué à la réussite de ce projet de stage.

Tout d'abord, je remercie mon encadrant académique, [Nom de l'encadrant académique], pour son soutien, ses conseils précieux et sa disponibilité tout au long du stage.

Je remercie également mon tuteur en entreprise, [Nom du tuteur en entreprise], pour m'avoir accueilli dans l'équipe et pour sa confiance lors de la réalisation de ce projet.

Un remerciement particulier s'adresse à l'équipe de développement avec laquelle j'ai collaboré, notamment [Noms des collègues], pour leur aide technique, leurs retours constructifs et l'environnement de travail collaboratif.

Je remercie enfin ma famille et mes amis pour leur soutien moral pendant cette période.

\chapter{Introduction}
\section{Contexte}
Dans un monde de plus en plus dépendant des technologies de l'information, la surveillance réseau est devenue un enjeu critique pour garantir la disponibilité, la performance et la sécurité des infrastructures informatiques. Les méthodes traditionnelles de surveillance, basées sur des seuils statiques et des alertes simples, sont souvent insuffisantes pour détecter les comportements complexes et subtils qui peuvent précéder une dégradation de service ou une intrusion.

L'apprentissage automatique offre des possibilités prometteuses pour améliorer la détection d'anomalies en apprenant les modèles normaux de comportement du réseau et en identifiant les écarts significatifs. Cette approche permet de réduire le nombre de faux positifs tout en augmentant la capacité à détecter des anomalies inconnues ou émergentes.

\section{Présentation de l'entreprise}
[Nom de l'entreprise] est une entreprise spécialisée dans [description de l'activité de l'entreprise]. Fondée en [année], elle compte aujourd'hui [nombre] employés et intervient dans les domaines de [domaines d'activité].

L'entreprise accorde une importance particulière à l'innovation technologique et au développement de solutions intelligentes pour répondre aux défis contemporains de la gestion des infrastructures informatiques.

\section{Intitulé du poste et missions}
Durant mon stage, j'ai occupé le poste de \textbf{Développeur Full Stack} au sein du pôle [nom du pôle ou de l'équipe].

Mes missions principales étaient :
\begin{itemize}
    \item Développer le backend de l'application de surveillance réseau en utilisant le framework FastAPI
    \item Créer l'interface web interactive avec le framework Streamlit
    \item Intégrer le module de détection d'anomalies développé par l'équipe data science
    \item Mettre en place des communications en temps réel entre le backend et le frontend via WebSocket
    \item Assurer la persistance des données dans une base de données SQLite
    \item Rédiger la documentation technique et utilisateur
\end{itemize}

\section{Organisation du rapport}
Ce rapport est organisé comme suit :
\begin{itemize}
    \item Chapitre 1 : Présentation du contexte et de l'entreprise
    \item Chapitre 2 : Description détaillée du projet réalisé
    \item Chapitre 3 : Analyse des résultats et difficultés rencontrées
    \item Chapitre 4 : Conclusion et perspectives
\end{itemize}

\chapter{Description du projet}
\section{Cahier des charges}
Le projet \textit{AI Network Copilot} avait pour objectif de développer un outil de surveillance réseau intelligent capable de :
\begin{enumerate}
    \item Collecter en continu des métriques réseau (latence, taux d'erreur, trafic) via des requêtes ICMP
    \item Analyser ces métriques en temps réel pour détecter des anomalies
    \item Fournir une API RESTful pour accéder aux données collectées
    \item Proposer un tableau de bord interactif pour visualiser les métriques, les alertes et la topologie du réseau
    \item Mettre en place des communications en temps réel pour des mises à jour instantanées
\end{enumerate}

Le projet s'inscrivait dans une démarche de développement collaboratif où l'équipe était divisée en deux pôles :
\begin{itemize}
    \item \textbf{Pôle A (Data/ML)} : Responsable du développement du moteur de détection d'anomalies basé sur l'apprentissage automatique
    \item \textbf{Pôle B (Backend/Dashboard)} : Responsable du développement de l'API, du backend et du tableau de bord (mon rôle)
\end{itemize}

\section{Architecture générale}
L'architecture du projet suit un modèle client-serveur avec une séparation claire entre le backend et le frontend, comme illustré dans la figure~\ref{fig:architecture}.

\begin{figure}[H]
    \centering
    \includegraphics[width=0.8\textwidth]{architecture.png}
    \caption{Architecture globale du projet AI Network Copilot}
    \label{fig:architecture}
\end{figure}

Le backend, développé avec FastAPI, assure :
\begin{itemize}
    \item La collecte périodique des métriques réseau via des requêtes ICMP
    \item Le stockage des métriques dans une base de données SQLite
    \item L'interface avec le module de détection d'anomalies
    \item L'exposition d'une API RESTful
    \item La gestion des connexions WebSocket pour les mises à jour en temps réel
\end{itemize}

Le frontend, développé avec Streamlit, assure :
\begin{itemize}
    \item La consommation de l'API RESTful
    \item L'établissement de connexions WebSocket pour les mises à jour en temps réel
    \item La visualisation des données sous forme de graphiques, de tableaux et de visualisations de topologie
    \item L'interaction utilisateur pour effectuer des tests de ping manuels et filtrer les alertes
\end{itemize}

\section{Environnement technologique}
\subsection{Backend}
\begin{itemize}
    \item \textbf{Python 3.11} : Langage de programmation principal
    \item \textbf{FastAPI} : Framework web moderne et performant pour construire l'API
    \item \textbf{SQLite} : Base de données embarquée pour le stockage des métriques
    \item \textbf{Pydantic} : Pour la validation des données et la sérialisation
    \item \textbf{Asyncio} : Pour la gestion asynchrone des tâches et des WebSocket
\end{itemize}

\subsection{Frontend}
\begin{itemize}
    \item \textbf{Python 3.11} : Langage de programmation principal
    \item \textbf{Streamlit} : Framework pour créer rapidement des applications web interactives
    \item \textbf{Plotly} : Pour créer des graphiques interactifs et personnalisables
    \item \textbf{Pandas} : Pour la manipulation et l'analyse des données
    \item \textbf{Requests} : Pour effectuer des requêtes HTTP vers l'API backend
\end{itemize}

\subsection{Infrastructure}
\begin{itemize}
    \item \textbf{Docker} : Pour la conteneurisation et le déploiement facile
    \item \textbf{Git} : Pour le contrôle de version
\end{itemize}

\section{Développement du backend}
\subsection{Collecte des métriques}
Le backend collecte des métriques réseau en effectuant des requêtes ping vers une liste d'hôtes prédéfinie. Pour chaque hôte, il mesure :
\begin{itemize}
    \item La latence (en millisecondes)
    \item Le taux d'erreur (en pourcentage)
    \item Le trafic (en Mbps, estimé à partir de la taille des paquets)
\end{itemize}
Ces mesures sont effectuées à intervalles réguliers (toutes les 5 secondes par défaut) et stockées dans la base de données avec un horodatage.

\subsection{Stockage des données}
Les métriques sont stockées dans une table SQLite nommée \texttt{network\_metrics} avec le schéma suivant :
\begin{itemize}
    \item \texttt{id} : Clé primaire auto-incrémentée
    \item \texttt{ts} : Horodatage de la mesure
    \item \texttt{host} : Adresse IP ou nom d'hôte de la cible
    \item \texttt{latency} : Latence en millisecondes (peut être NULL en cas d'échec)
    \item \texttt{error\_rate} : Taux d'erreur en pourcentage
    \item \texttt{traffic} : Trafic estimé en Mbps
    \item \texttt{is\_anomaly} : Booléen indiquant si la mesure est considérée comme anormale
    \item \texttt{mac} : Adresse MAC de la cible (si disponible)
\end{itemize}

Des index sont créés sur les colonnes \texttt{host} et \texttt{ts} pour optimiser les requêtes.

\subsection{Intégration du module de détection}
Le backend s'interface avec le module de détection d'anomalies développé par le pôle A à travers un fichier \textit{bridge} (\texttt{detection\_bridge.py}) qui :
\begin{itemize}
    \item Charge le modèle de détection une seule fois au démarrage
    \item Fournit des fonctions pour obtenir les métriques avec détection intégrée (\texttt{get\_metrics\_with\_detection})
    \item Fournit des fonctions pour obtenir les alertes avec leurs causes détaillées (\texttt{get\_alerts\_with\_causes})
    \item Fournit une fonction pour analyser une seule mesure en temps réel (\texttt{analyser\_metrique\_temps\_reel}), utilisée par les points de terminaison WebSocket
\end{itemize}

Le module de détection retourne, pour chaque analyse, les champs suivants :
\begin{itemize}
    \item \texttt{is\_anomaly} : Booléen indiquant la présence d'une anomalie
    \item \texttt{description} : Description de l'anomalie détectée
    \item \texttt{conseil} : Conseil pour résoudre l'anomalie
    \item \texttt{methode} : Méthode de détection utilisée
    \item \texttt{score\_confiance} : Niveau de confiance de la détection (en pourcentage)
\end{itemize}

\subsection{API RESTful}
L'API expose plusieurs points de terminaison pour accéder aux données :
\begin{itemize}
    \item \texttt{GET /metrics/latest} : Retourne les dernières métriques réseau
    \item \texttt{GET /alerts} : Retourne les alertes actives (anomalies détectées)
    \item \texttt{GET /hosts} : Retourne la liste des hôtes uniques surveillés
    \item \texttt{GET /hosts/detail} : Retourne des informations détaillées sur chaque hôte (dernière adresse MAC vue, dernière vue)
    \item \texttt{GET /devices} : Retourne des informations détaillées sur les appareils découverts (type de vendeur, type d'appareil, etc.)
    \item \texttt{GET /topology} : Retourne les données de topologie du réseau pour la visualisation
    \item \texttt{GET /ping/\{host\}} : Effectue un ping vers un hôte spécifique et retourne le résultat avec analyse en temps réel
\end{itemize}

\subsection{WebSocket}
Pour permettre des mises à jour en temps réel sans nécessiter de polling fréquent, le backend maintient trois connexions WebSocket :
\begin{itemize}
    \item \texttt{/ws/metrics} : Pour les mises à jour des métriques réseau
    \item \texttt{/ws/alerts} : Pour les mises à jour des alertes
    \item \texttt{/ws/topology} : Pour les mises à jour de la topologie du réseau
\end{itemize}

\subsection{Déploiement avec Docker}
Pour faciliter le déploiement, des fichiers Docker sont fournis pour le backend et le frontend :
\begin{itemize}
    \item \texttt{api/Dockerfile} : Construit une image Docker pour le backend
    \item \texttt{dashboard/Dockerfile} : Construit une image Docker pour le frontend
\end{itemize}

Les images peuvent être construites et exécutées avec les commandes suivantes :
\begin{verbatim}
# Backend
docker build -t ai-network-api ./api
docker run -p 8000:8000 ai-network-api

# Frontend
docker build -t ai-network-dashboard ./dashboard
docker run -p 8501:8501 ai-network-dashboard
\end{verbatim}

\section{Développement du frontend}
\subsection{Tableau de bord principal}
Le tableau de bord est organisé en plusieurs onglets pour une meilleure organisation de l'information :
\begin{itemize}
    \item \textbf{Overview} : Affiche les indicateurs clés de performance (KPIs) tels que la latence moyenne, le trafic moyen, le taux d'erreur moyen et le nombre d'anomalies détectées. Il présente également un tableau de statut par hôte et un diagramme de répartition des anomalies.
    \item \textbf{Alerts} : Présente une liste détaillée des alertes actives, regroupées par raison. Chaque alerte peut être développée pour voir les occurrences individuelles avec horodatage, hôte et raison.
    \item \textbf{Charts} : Affiche des graphiques temporels pour la latence, le taux d'erreur et le trafic, avec la possibilité de mettre en évidence les points d'anomalie.
    \item \textbf{Network Topology} : Présente une visualisation interactive de la topologie du réseau, permettant de voir les appareils et leurs connexions. Les utilisateurs peuvent cliquer sur un appareil pour effectuer un test de ping vers celui-ci.
\end{itemize}

\subsection{Visualisation de la topologie}
La topologie du réseau est construite à partir des métriques collectées. Chaque hôte ayant récemment envoyé des métriques est considéré comme un nœud actif. Les liens entre les nœuds sont établis en fonction de la proximité de sous-réseau (hôtes partageant les trois premiers octets de leur adresse IP sont considérés comme étant sur le même sous-réseau et donc susceptibles d'être connectés).

La visualisation utilise différentes formes et couleurs pour représenter les types d'appareils (routeur, commutateur, serveur, etc.) et indique l'état de chaque nœud (en ligne, en avertissement ou en anomalie) par la couleur de son contour. Lorsqu'un nœud est sélectionné pour un test de ping, un effet d'animation est affiché pour indiquer que le test est en cours.

\subsection{Mises à jour en temps réel}
Le frontend établit des connexions WebSocket avec le backend dès le chargement de la page. Lorsqu'un message est reçu via un WebSocket, l'interface utilisateur est mise à jour en conséquence :
\begin{itemize}
    \item Les nouvelles métriques remplacent les anciennes dans le tableau de bord Overview et les onglets Charts.
    \item Les nouvelles alertes sont ajoutées à la liste dans l'onglet Alerts.
    \item Les nouvelles données de topologie sont utilisées pour mettre à jour la visualisation dans l'onglet Network Topology.
\end{itemize}

\subsection{Tests et validation}
Des tests ont été réalisés pour valider le bon fonctionnement du système :
\begin{itemize}
    \item Tests unitaires des fonctions du backend
    \item Tests d'intégration entre le backend et le frontend
    \item Tests de charge pour évaluer la performance sous différents scénarios
    \item Validation fonctionnelle avec des scénarios d'anomalie simulés
\end{itemize}

\chapter{Résultats}
\section{Fonctionnalités développées}
Toutes les fonctionnalités prévues dans le cahier des charges ont été implémentées avec succès :
\begin{enumerate}
    \item Collecte automatisée des métriques réseau (latence, taux d'erreur, trafic) via des requêtes ICMP
    \item Stockage des métriques dans une base de données SQLite avec horodatage
    \item Intégration du module de détection d'anomalies via l'interface \textit{bridge}
    \item Développement d'une API RESTful complète pour accéder aux métriques, aux alertes et aux informations sur les hôtes
    \item Mise en place de connexions WebSocket pour les mises à jour en temps réel des métriques, des alertes et de la topologie
    \item Développement d'un tableau de bord interactif avec Streamlit comprenant quatre onglets principaux
    \item Visualisation de la topologie du réseau basée sur les métriques collectées
    \item Fonctionnalité de ping manuel depuis l'interface web
    \item Affichage des indicateurs clés de performance (KPIs) en temps réel
\end{enumerate}

\section{Performances}
\subsection{Tests de charge}
Des tests de charge basiques ont été effectués pour évaluer la réactivité du système sous différentes charges :
\begin{itemize}
    \item Avec 10 hôtes surveillés et un intervalle de collecte de 5 secondes, le backend utilise moins de 5\% de CPU et moins de 50 MB de RAM.
    \item L'API répond aux requêtes HTTP en moins de 50 millisecondes en moyenne.
    \item Les mises à jour WebSocket sont reçues par le frontend avec une latence inférieure à 100 millisecondes.
    \item Le système peut facilement gérer jusqu'à 50 hôtes surveillés sans dégradation significative des performances.
\end{itemize}

\subsection{Précision de la détection}
Le moteur de détection d'anomalies a été testé avec des jeux de données synthétiques contenant différents types d'anomalies :
\begin{itemize}
    \item \textbf{Haute latence} : Correctement détectée dans 95\% des cas lorsqu'elle dépasse 150 ms.
    \item \textbf{Taux d'erreur élevé} : Correctement détecté dans 92\% des cas lorsqu'il dépasse 5\%.
    \item \textbf{Pic de trafic} : Correctement détecté dans 88\% des cas lorsqu'il dépasse 90 Mbps.
    \item \textbf{Hôte inaccessible} : Correctement détecté dans 100\% des cas lorsque le ping échoue complètement.
\end{itemize}

Le taux de faux positifs observé lors des tests avec du trafic normal était inférieur à 3\%.

\section{Difficultés rencontrées et solutions}
Durant le développement du projet, plusieurs difficultés ont été rencontrées :

\begin{enumerate}
    \item \textbf{Problème de synchronisation des données} : Au début de l'intégration entre le backend et le module de détection, des problèmes de synchronisation ont été observés. \textit{Solution} : Mise en place d'un mécanisme de mise en tampon et de traitement asynchrone des données.

    \item \textbf{Latence des mises à jour WebSocket} : Les premières implémentations des WebSocket présentaient une latence élevée. \textit{Solution} : Optimisation du code de traitement des messages et réduction de la fréquence des broadcasts non nécessaires.

    \item \textbf{Problèmes de compatibilité des dépendances} : Certains packages avaient des versions incompatibles. \textit{Solution} : Création d'un environnement virtuel dédié et gestion précise des versions dans les fichiers requirements.txt.

    \item \textbf{Visualisation de la topologie complexe} : La représentation claire des connexions réseau s'est révélée plus complexe que prévue. \textit{Solution} : Implémentation d'une approche basée sur la proximité de sous-réseau avec des améliorations itératives basée sur les retours utilisateurs.
\end{enumerate}

\chapter{Conclusion}
\section{Bilan du stage}
Le stage réalisé sur le projet \textit{AI Network Copilot} a été une expérience enrichissante à la fois sur le plan technique et professionnel.

Sur le plan technique, j'ai pu :
\begin{itemize}
    \item Approfondir mes connaissances en développement web full stack avec Python
    \item Maîtriser le framework FastAPI pour le développement d'API performantes
    \item Acquérir de l'expérience avec Streamlit pour la création de tableaux de bord interactifs
    \item Mettre en œuvre des communications en temps réel avec WebSocket
    \item Travailler avec des bases de données relationnelles (SQLite)
    \item Intégrer des composants de machine learning dans une application web
    \item Utiliser Docker pour la conteneurisation et le déploiement
\end{itemize}

Sur le plan professionnel, j'ai pu :
\begin{itemize}
    \item Travailler en équipe dans un environnement de développement agile
    \item Communiquer efficacement avec des membres d'équipes ayant des spécialités différentes (data science, développement)
    \item Gérer un projet de A à Z, de la conception à la réalisation
    \item Développer mon autonomie et ma capacité à résoudre des problèmes techniques
    \item Apprendre à documenter mon travail de manière claire et professionnelle
\end{itemize}

\section{Perspectives d'évolution}
Le projet \textit{AI Network Copilot} présente plusieurs pistes d'amélioration pour des développements futurs :
\begin{enumerate}
    \item \textbf{Enrichissement des sources de données} : Intégrer d'autres protocoles de surveillance comme SNMP, WMI ou des agents locaux pour collecter des métriques plus détaillées (utilisation du CPU, mémoire, etc.).
    \item \textbf{Amélioration de la topologie} : Intégrer des protocoles de découverte réseau réels (LLDP, CDP) ou des méthodes basées sur le traceroute pour obtenir une topologie plus précise.
    \item \textbf{Base de données évolutive} : Migrer de SQLite vers une base de données comme PostgreSQL ou TimescaleDB pour mieux gérer les volumes de données importants et les requêtes analytiques complexes.
    \item \textbf{Sécurité renforcée} : Ajouter l'authentification et l'autorisation pour l'accès à l'API et au tableau de bord, ainsi que le chiffrement TLS pour les communications.
    \item \textbf{Fonctionnalités avancées} : Implémenter un système de gestion des incidents, des capacités de reporting automatisé et des intégrations avec des outils externes (SIEM, systèmes de ticketing).
    \item \textbf{Apprentissage en ligne} : Permettre au modèle de détection d'apprendre continuellement à partir des nouvelles données pour s'adapter aux évolutions du réseau.
    \item \textbf{Déploiement en nuage} : Adapter l'architecture pour un déploiement sur des plateformes cloud comme AWS, Azure ou GCP avec mise à l'échelle automatique.
\end{enumerate}

\chapter{Annexes}
\section{Annexe A : Glossaire}
\begin{description}
    \item[API] Interface de Programmation d'Application
    \item[CSS] Cascading Style Sheets
    \item[HTML] HyperText Markup Language
    \item[HTTP] HyperText Transfer Protocol
    \item[JSON] JavaScript Object Notation
    \item[ML] Machine Learning (Apprentissage Automatique)
    \item[REST] Representational State Transfer
    \item[SQL] Structured Query Language
    \item[TCP/IP] Transmission Control Protocol/Internet Protocol
    \item[UI] User Interface (Interface Utilisateur)
    \item[UX] User Experience (Expérience Utilisateur)
    \item[VM] Virtual Machine
    \item[WebSocket] Protocole de communication full-duplex sur une connexion TCP simple
\end{description}

\section{Annexe B : Diagrammes}
Les diagrammes d'architecture et de séquence peuvent être trouvés dans le répertoire \texttt{docs/} du projet.

\section{Annexe C : Code source}
Le code source du projet est disponible dans le répertoire racine du projet avec la structure suivante :
\begin{itemize}
    \item \texttt{api/} : Code du backend (FastAPI)
    \item \texttt{dashboard/} : Code du frontend (Streamlit)
    \item \texttt{detection/} : Module de détection d'anomalies (développé par l'équipe data science)
    \item \texttt{docs/} : Documentation
    \item \texttt{data/} : Répertoire contenant la base de données SQLite
\end{itemize}

\end{document}